\lecture{33}{Wed 30 Nov 2022 11:30}{Series}

\section{Series}

\subsection{Definitions}

\begin{definition}(Series of Numbers)
	Define the \textit{partial sum} \[
		S_n = \sum_{k=1}^{n} x_k
	.\] 
	We say \[
		S = \sum_{k=1}^{\infty} x_k :\iff S_n \to S \text{ as } n\to \infty 
	.\] 
\end{definition}

\begin{definition}[Absolutely Convergent]
	We say a series  $\sum_{k=1}^{\infty} x_k $ is \textit{absolutely convergent} if $\sum_{k=1}^{\infty} \left| x_k \right|  $.
\end{definition}

\begin{lemma}
	Let $\sum_{k=1}^{\infty} x_n$ be a series with $x_n \ge 0$. Then it is convergent iff there exists $M$ such that $S_n \le M$ for all $n$. Then $\sum_{k=1}^{\infty} x_n \le M$.
\end{lemma}
\begin{proof}
	By Monotone Convergence Theorem applied to $S_n$.
\end{proof}

\begin{lemma}[Cauchy Criterion]
	$\sum x_n$ converges iff for any $\varepsilon>0$ ther is $M_\varepsilon$ such that $\sum_{k=n}^m x_k$ if $n\ge m\ge M_\varepsilon$.
\end{lemma}
\begin{proof}
	$S_n$ converges iff $S_n$ is cauchy sequence.
\end{proof}

\begin{proposition}
	Absolute convergnce implies convergence.
\end{proposition}
\begin{proof}
	Use triangle inequality, prove that $\sum x_n$ is Cauchy.
\end{proof}

\begin{definition}[Rearrangement of Series]
	Take $\varphi: \mathbb{N}\to \mathbb{N}$ be a bijection. The series
	\[
		\sum_{n=1}^{\infty} x_{\varphi(n)}
	\] is called a \textit{rearrangement} of $\sum_{n=1}^{\infty} x_n $.
\end{definition}

\begin{theorem}
	If $x_n \ge 0$, $\sum_{n=1}^{\infty} x_n < \infty $, then $\sum_{n=1}^{\infty} x_{\varphi(n)} < \infty $ and $\sum_{n=1}^{\infty} x_{n} = \sum_{n=1}^{\infty} x_{\varphi(n)}$
\end{theorem}
\begin{proof}
	Let $S_n' = X_{\varphi(1)} + \ldots + X_{\varphi(n)}$. Let $N = \max \{\varphi(1), \ldots, \varphi(n) \} $.
	Hence for any $n$,
	\[
		S_n' \le S_N \le M
	.\] 
	Thus $\sum_{n=1}^{\infty} x_{\varphi(n)} = \lim_{n \to \infty} S_n' \le M$.
\end{proof}
\begin{theorem}
	If $\varphi: \mathbb{N}\to \mathbb{N}$ is bijection, then $\sum x_n$ absolutely convergent iff $\sum x_{\varphi(n)}$ absolutely convergent, and then 
	\[
		\sum_{n=1}^{\infty} x_n 
		= \sum_{n=1}^{\infty} x_{\varphi(n)}
	.\] 
\end{theorem}
\begin{proof}
	Convergence follows from the proposition above. To prove the sums are equal, split into positive and negative parts. 
\end{proof}

\begin{definition}[Conditionally Convergent]
	$\sum x_n$ is convergent, but not absolutely.
\end{definition}
Equivalently, both $\sum_{n=1}^{\infty } x_n^+$ and $\sum_{n=1}^{\infty} x_n^-$ are infinite.

\begin{theorem}
	If $\sum_{n=1}^{\infty} x_n$ converges conditionally, then for each real number $x$ there exists $\varphi: \mathbb{N}\to \mathbb{N}$ such that
	\[
		\sum_{n=1}^{\infty} x_{\varphi_x(n)} = x
	.\] 
\end{theorem}

\begin{proposition}
	$\sum a_n$ convergent implies $a_n \to 0$.
\end{proposition}
\begin{proof}
	$a_n = S_n - S_{n-1} \to S - S = 0$
\end{proof}

\begin{example}
	\begin{enumerate}
	\item (Geometric series)
	\item ($p$-series)
		\[
		\sum_{n=1}^{\infty} \frac{1}{n^p}
		.\] 
		If $p\le 0$, then $\frac{1}{n^p}\ge 1 \not \to 0$ divergent.
		If $p>0$, $\frac{1}{n^p} \to 0$. But is the series it convergent?

		Apply \textit{Cauhy's condensation method}: 
		\[
		S_{2^k}\ge 1 + \frac{1}{2} \left( \frac{1}{2^{p-1} + \frac{1}{(2^{p-1})^2}} + \ldots + \frac{1}{(2^{p-1})^k} \right) 		.\] 
		 \[
			S_{2^k} \ge  1 + \frac{1}{2} \left( a + a^2 + \ldots + a^k \right) \qquad a = \frac{1}{2^{p-1}}
		.\] 
		So the series diverges if $p \le 1$.
	\end{enumerate}
\end{example}

\begin{theorem}[Cauchy's Condensation Test]
	Let $a_n \ge 0$, $a_n$ strictly decreasing. Then $\sum a_n <\infty $ iff $\sum 2^k a_{2^k} < \infty $.
	
\end{theorem}

\begin{theorem}[Comparison Theorem]
	Suppose $\sum x_n$, $\sum y_n$, $x_n, y_n \ge 0$ are such that  $x_n \le y_n$ for $n \ge N$. Then 
	\[
	\sum y_n < \infty \implies \sum x_n < \infty 
	.\] 
	Also
	\[
		\sum x_n = \infty  \implies \sum y_n = \infty
	.\] 
\end{theorem}
\begin{proof}
	work with partial sums.
\end{proof}


\begin{theorem}[Limit Comparison Theorem]
	Suppose $\limsup \frac{x_n}{y_n}< \infty$, then
	\[
		\sum y_n <\infty  \implies \sum x_n < \infty
	.\] 
	SUppose $\liminf \frac{x_n}{y_n}>0$ then
	\[
		\sum y_n = \infty  \implies \sum x_n = \infty 
	.\] 
\end{theorem}
\begin{theorem}[Root Test]
	Let $x_n \ge 0$ and 
	\[
		\rho =\limsup _{n \to  \infty } |x_n|^{\frac{1}{n}}
	.\] 
	Then
	\begin{enumerate}
		\item If $\rho < 1$ then $\sum x_n $ absolutely convergent
		\item If  $\rho > 1$ then $\sum x_n$ divergent.
		\item Otherwise the test is conclusive.
	\end{enumerate}
\end{theorem}
\begin{proof}
	\begin{enumerate}
		\item Choose $\rho < \tilde \rho < 1$. Then there exists $N $ such that for $n \ge N$,
	\[
		|x_n|^{\frac{1}{n}}\le \rho \tilde
	.\] 
	Hence $|x_n|\le \tilde \rho^n$. Now we can apply comparison test.
\item Pick subsequence such that $|x_{n_k}|^{\frac{1}{n}} \ge 1$, so $x_n \not \to 0$.
	\end{enumerate}
\end{proof}

